\documentclass[11pt]{article}
\usepackage{fontspec}
\setmainfont{Latin Modern Roman}
\usepackage{amsmath,amssymb,mathtools,unicode-math}
\setmathfont{Latin Modern Math}
\usepackage{xurl,hyperref}
\hypersetup{hidelinks,pdftitle={Problem statement}}
\newcommand{\sref}[2]{\hyperlink{source:#1:#2}{[S#2]}}
\newcommand{\cataloguescope}[1]{\paragraph{Mathematical question.}#1}
\begin{document}
\section{Goldberg's conjecture}\label{statement}

\subsection{Definitions and mathematical statement}

\paragraph{Shared notation.}$\mathbb N=\{1,2,\ldots\}$, and $\mathbb Z,\mathbb Q,\mathbb R,\mathbb C$ denote the integers, rationals, reals and complex numbers. Any other conventions are specified locally.


For a finite loopless multigraph $G$ of maximum degree $\Delta(G)$ and chromatic index $\chi'(G)$, define \[w(G)=\max_{H\subseteq G,\ |V(H)|\ge2}\left\lceil\frac{|E(H)|}{\lfloor |V(H)|/2\rfloor}\right\rceil,\] taking $w(G)=0$ if no such subgraph exists. Restricting to at least two vertices makes the displayed denominator nonzero; even-order subgraphs give a value at most $\Delta(G)$.
\paragraph{Statement recorded in the supplied source.}
The main recorded conjecture is \[\chi'(G)\le\max\{\Delta(G)+1,w(G)\}.\] The same source gives two related conjectural bounds, \[\chi'(G)\le\max\{\Delta(G),w(G)\}+1\quad\text{and}\quad\chi'(G)\le\max\{\Delta(G),w(G)+1\}.\] \sref{statement}{2}
\subsection{Short English statement}
Does maximum degree together with the densest subgraph obstruction determine the proposed edge-coloring bounds for multigraphs?
\subsection{Sources}
\begin{description}
\item[\hypertarget{source:statement:1}{S1}] ULAM.AI and the UnsolvedMath Contributors, UnsolvedMath: A Curated Collection of Open Mathematics Problems, catalogue entry. CC BY 4.0. \url{https://huggingface.co/datasets/ulamai/UnsolvedMath}.
\item[\hypertarget{source:statement:2}{S2}] Open Problem Garden, Goldberg's conjecture, original node 2242, statement and mathematical discussion. \url{https://www.openproblemgarden.org/op/goldbergs_conjecture}.
\end{description}
\label{end:statement}
\end{document}
